\documentclass[11pt]{article}

\usepackage[T1]{fontenc}
\usepackage{amsmath,amssymb,amsthm,mathtools}
\usepackage{geometry}
\usepackage{graphicx}
\usepackage{microtype}
\usepackage[hidelinks]{hyperref}
\geometry{margin=1in}

\newtheorem{theorem}{Theorem}
\newtheorem{lemma}{Lemma}
\theoremstyle{remark}
\newtheorem{remark}{Remark}

\newcommand{\C}{\mathbb C}
\newcommand{\cS}{\mathcal S}
\newcommand{\id}{\mathrm{Id}}
\newcommand{\tr}{\operatorname{Tr}}

\title{The quadratic quantum Wasserstein channel complexity\\
is not convex}
\author{Candidate counterexample for arXiv:2402.13150\\
Run \texttt{fa\_banach\_001}, lane 10, model GPT5.6}
\date{22 July 2026}

\begin{document}
\maketitle

\begin{center}
\fbox{\parbox{0.92\linewidth}{
\textbf{Status.} Candidate counterexample likely valid; full negative answer
to the channel-complexity convexity question in Section~5.1 of
arXiv:2402.13150.  On one qubit, for the source paper's symmetric Pauli cost,
the mixture of the identity channel and the Pauli-$X$ channel violates the
convexity inequality.  All estimates are exact and finite-dimensional.
}}
\end{center}

\section{Source question}

Bunth, Pitrik, Titkos, and Virosztek define, for a finite collection of
observables $\mathcal A=\{A_1,\ldots,A_k\}$, the channel Wasserstein
complexity
\begin{equation}\label{eq:complexity}
 C_W(\Phi)=\max_{\rho\in\cS(\mathcal H)}
 d_{\mathcal A}(\rho,\Phi(\rho)),
\end{equation}
where $d_{\mathcal A}$ is their modified quadratic quantum Wasserstein
divergence \cite{BPVT}.  At the end of Section~5.1 (printed page 18), they
write:
\begin{quote}
``We note that whether the Wasserstein complexity (37) is convex or not is an
open question.''
\end{quote}
The original passage is reproduced in Figure~\ref{fig:problem}.

\begin{figure}[ht]
  \centering
  \includegraphics[width=0.98\linewidth]{figures/open_problem_crop.png}
  \caption{The open question in arXiv:2402.13150v4, printed page 18.}
  \label{fig:problem}
\end{figure}

Here convexity has its standard meaning on the convex set of quantum
channels:
\[
 C_W((1-t)\Phi_0+t\Phi_1)
 \leq (1-t)C_W(\Phi_0)+tC_W(\Phi_1),\qquad 0\leq t\leq1.
\]

\section{Definitions and proof intuition}

For two states $\rho,\omega$, let
\[
 D_{\mathcal A}^2(\rho,\omega)
 =\min_{\Pi\in\mathcal C(\rho,\omega)}\tr(C_{\mathcal A}\Pi),
 \qquad
 C_{\mathcal A}=\sum_j(A_j\otimes I-I\otimes A_j^T)^2,
\]
with the coupling convention of \cite{BPVT}.  The modified divergence is
\begin{equation}\label{eq:ddef}
 d_{\mathcal A}^2(\rho,\omega)
 =D_{\mathcal A}^2(\rho,\omega)
 -\frac12\bigl(D_{\mathcal A}^2(\rho,\rho)
                 +D_{\mathcal A}^2(\omega,\omega)\bigr).
\end{equation}

\paragraph{Proof intuition.}
Take the source paper's own symmetric qubit cost
$\mathcal A_s=\{\sigma_x,\sigma_y,\sigma_z\}$.  Mix the identity channel with
the bit flip $\mathsf X(\rho)=\sigma_x\rho\sigma_x$.  Convexity would make the
complexity of this mixture grow at most linearly in the mixing parameter
$t$.  But on the pure input $|0\rangle\langle0|$, the output immediately
acquires a second eigenvalue $t$.  The square root in the definition of
$d_s$ produces an exact term $\sqrt{t(1-t)}$ in $d_s^2$, and hence
$d_s\asymp t^{1/4}$.  This is much larger than a linear function near zero.

Only one global estimate is needed: the Pauli cost operator equals
$8(I-|\Omega\rangle\langle\Omega|)$, so every channel has complexity at most
$2\sqrt2$.  At $t=1/4$, this upper bound on the convexity right-hand side is
$\sqrt2/2$, while the single pure input gives the exact lower bound
$\sqrt{1+\sqrt3}$.

\section{Two elementary Pauli calculations}

Write
\[
 \sigma_x=\begin{pmatrix}0&1\\1&0\end{pmatrix},\quad
 \sigma_y=\begin{pmatrix}0&-i\\i&0\end{pmatrix},\quad
 \sigma_z=\begin{pmatrix}1&0\\0&-1\end{pmatrix},
\]
and abbreviate $D_s=D_{\mathcal A_s}$ and $d_s=d_{\mathcal A_s}$ as in the
source paper.

\begin{lemma}[Pure-input formula]\label{lem:pure}
If at least one of $\rho,\omega\in\cS(\C^2)$ is pure, then
\begin{equation}\label{eq:pure}
 d_s^2(\rho,\omega)=\sum_{j=x,y,z}
 \left\{
 \tr(\sqrt\rho\,\sigma_j\sqrt\rho\,\sigma_j)
 +\tr(\sqrt\omega\,\sigma_j\sqrt\omega\,\sigma_j)
 -2\tr(\rho\sigma_j)\tr(\omega\sigma_j)
 \right\}.
\end{equation}
\end{lemma}

\begin{proof}
When one marginal is pure, its only coupling with the other state is the
product coupling.  Expanding its cost gives
\[
 D_s^2(\rho,\omega)=\sum_j
 \bigl(\tr(\rho\sigma_j^2)+\tr(\omega\sigma_j^2)
 -2\tr(\rho\sigma_j)\tr(\omega\sigma_j)\bigr).
\]
For any state $\tau$, the canonical purification realizes the self-cost and
gives
\[
 D_s^2(\tau,\tau)=2\sum_j
 \bigl(\tr(\tau\sigma_j^2)
 -\tr(\sqrt\tau\,\sigma_j\sqrt\tau\,\sigma_j)\bigr).
\]
Substitution in \eqref{eq:ddef} proves \eqref{eq:pure}.  These are also the
source paper's displayed pure-state and self-cost formulas
\cite[Sec.~2]{BPVT}.
\end{proof}

\begin{lemma}[Uniform diameter bound]\label{lem:diameter}
For every pair of qubit states, $d_s(\rho,\omega)\leq2\sqrt2$.  Consequently,
$C_W(\Phi)\leq2\sqrt2$ for every qubit channel $\Phi$.
\end{lemma}

\begin{proof}
Let
\[
 |\Omega\rangle=\frac1{\sqrt2}(|00\rangle+|11\rangle).
\]
The Pauli expansion of the maximally entangled projection is
\[
 \sum_{j=x,y,z}\sigma_j\otimes\sigma_j^T
 =4|\Omega\rangle\langle\Omega|-I.
\]
Since $\sigma_j^2=I$,
\begin{align*}
 C_s&=\sum_j(\sigma_j\otimes I-I\otimes\sigma_j^T)^2\\
 &=6I-2\sum_j\sigma_j\otimes\sigma_j^T
 =8(I-|\Omega\rangle\langle\Omega|).
\end{align*}
Thus $0\leq C_s\leq8I$.  Every coupling has cost at most $8$, so
$D_s^2(\rho,\omega)\leq8$.  Formula \eqref{eq:ddef} subtracts nonnegative
self-costs; hence $d_s^2\leq D_s^2\leq8$.  Maximizing over $\rho$ proves the
channel bound.  The same spectral decomposition appears in
\cite[Eq.~(24)]{BPVT}.
\end{proof}

\section{The counterexample}

Let $\mathsf X$ be the Pauli bit-flip channel and define
\begin{equation}\label{eq:channel}
 \mathsf X(\rho)=\sigma_x\rho\sigma_x,
 \qquad
 \Lambda_t=(1-t)\id+t\mathsf X.
\end{equation}

\begin{theorem}\label{thm:counterexample}
For the symmetric Pauli quadratic Wasserstein divergence on one qubit, the
channel complexity \eqref{eq:complexity} is not convex.  Explicitly,
\[
 C_W\!\left(\frac34\id+\frac14\mathsf X\right)
 \geq\sqrt{1+\sqrt3}
 >\frac{\sqrt2}{2}
 \geq\frac34C_W(\id)+\frac14C_W(\mathsf X).
\]
\end{theorem}

\begin{proof}
Use the pure input
\[
 \rho_0=|0\rangle\langle0|,
 \qquad
 \omega_t=\Lambda_t(\rho_0)
 =\begin{pmatrix}1-t&0\\0&t\end{pmatrix}.
\]
Apply Lemma~\ref{lem:pure}.  For $\sigma_z$, the three terms in braces are
\[
 1+1-2(1)(1-2t)=4t.
\]
For each of $\sigma_x$ and $\sigma_y$, both expectation values and the
$\rho_0$ square-root term vanish, while
\[
 \tr(\sqrt{\omega_t}\,\sigma_j\sqrt{\omega_t}\,\sigma_j)
 =2\sqrt{t(1-t)}.
\]
Therefore
\begin{equation}\label{eq:exactcurve}
 d_s^2(\rho_0,\Lambda_t(\rho_0))
 =4t+4\sqrt{t(1-t)},
\qquad
 C_W(\Lambda_t)\geq
 \sqrt{4t+4\sqrt{t(1-t)}}.
\end{equation}
At $t=1/4$, this lower bound is $\sqrt{1+\sqrt3}$.

On the other hand, $C_W(\id)=0$ because $d_s(\rho,\rho)=0$.  By
Lemma~\ref{lem:diameter}, $C_W(\mathsf X)\leq2\sqrt2$.  If $C_W$ were convex,
then
\[
 C_W(\Lambda_{1/4})
 \leq\frac34C_W(\id)+\frac14C_W(\mathsf X)
 \leq\frac{\sqrt2}{2}.
\]
This contradicts \eqref{eq:exactcurve}, since
$1+\sqrt3>1/2$.  Hence $C_W$ is not convex.
\end{proof}

\begin{remark}[Robust mechanism]
Equation \eqref{eq:exactcurve} gives
$C_W(\Lambda_t)\geq2t^{1/4}(1+o(1))$ as $t\downarrow0$, whereas convexity
would give $C_W(\Lambda_t)\leq2\sqrt2\,t$.  Thus the violation holds for an
interval of small positive $t$ and is not a numerical coincidence at
$t=1/4$.
\end{remark}

\section{Verification, scope, and novelty}

The proof is exact; no numerical optimization over couplings or states is
used.  The accompanying script \texttt{code/verify\_counterexample.py}
symbolically checks the Pauli cost identity, its spectrum, the three terms in
\eqref{eq:pure}, and the strict convexity gap at $t=1/4$.  It is a redundant
finite-dimensional check, not a substitute for the proof.

The counterexample answers exactly the question as stated: the observable
collection is the source paper's standard symmetric Pauli choice, both
endpoints and their mixture are quantum channels, and the same divergence
and maximization as equation (37) are used.  It does not assert that every
other quantum Wasserstein channel complexity is nonconvex, nor does it address
convexity notions for the different order-one/Hamming-Wasserstein complexity.

A bounded literature search was performed on 22 July 2026 using the exact
source sentence, the source title and arXiv id, and combinations of
``quantum Wasserstein complexity,'' ``convex,'' ``channel,'' ``Pauli,'' and
``quadratic divergence.''  The local run indexes were also searched.  No
prior answer to this exact Section~5.1 question was found.  Search hits about
convexity concerned different order-one or resource-dependent channel
complexities.  Novelty is therefore plausible but not certified.

The source paper's separate De Palma--Trevisan triangle-inequality conjecture
is no longer open: Wirth subsequently proved it in full generality
\cite{Wirth}.  That literature update is why the present packet targets the
still-distinct channel-complexity question rather than claiming a new proof of
the paper's main conjecture.

\paragraph{Human-review recommendation.}
High priority and low expected effort.  Check (i) the pure-input formula's
normalization against equation (37), (ii) the identity
$C_s=8(I-|\Omega\rangle\langle\Omega|)$, and (iii) that convexity is meant on
the affine space of channels.  These three checks reproduce the entire
argument.

\begin{thebibliography}{9}
\setlength{\itemsep}{0.6em}

\bibitem{BPVT}
G.~Bunth, J.~Pitrik, T.~Titkos, and D.~Virosztek,
\emph{On the metric property of quantum Wasserstein divergences},
arXiv:2402.13150v4 (2025), especially Sections 2, 3.1, and 5.1.

\bibitem{Wirth}
M.~Wirth,
\emph{Triangle Inequality for a Quantum Wasserstein Divergence},
arXiv:2511.20450 (2025).

\end{thebibliography}

\end{document}
